\documentclass[12pt]{article}
\usepackage[margin=1in]{geometry}
\usepackage{setspace}
\usepackage{booktabs}
\usepackage{amsmath}
\usepackage{amssymb}
\usepackage{natbib}
\usepackage{graphicx}
\usepackage{threeparttable}
\usepackage{hyperref}

\onehalfspacing
\setcitestyle{authoryear,open={(},close={)}}

\title{Institutional Choice in Postwar Japan: Welfare Loans versus Public Pawnshops, 1961}
\author{}
\date{}

\begin{document}

\maketitle

\begin{abstract}
How did the formation of the welfare state transform credit provision to low-income households? This paper examines institutional allocation in postwar Japan by analyzing a 1961 survey of 6,152 low-income households in Kanagawa prefecture. The central finding is that welfare loans and public pawnshops served observably different populations, selected by economic position and creditworthiness rather than shock type alone. Welfare loans were concentrated among households experiencing business failure or illness but possessing material foundations (mean assets: 3.5 items) and no public-assistance history; public pawnshops served asset-poor households (mean assets: 2.6 items) with precarious economic positions. Contrary to displacement narratives, welfare-loan users continued to layer multiple credit sources. Previous pawnshop users were 37.7 percentage points more likely to report ongoing pawning, while welfare-loan users retained informal borrowing networks. These patterns suggest that postwar financial modernization involved not the unification of credit institutions but rather their hierarchical stratification, with the state as arbiter of creditworthiness within the poor population.

\vspace{0.5cm}
\noindent \textbf{Keywords:} welfare state, household finance, credit allocation, Japan, economic history, institutional selection
\end{abstract}

\newpage
\section{Introduction}

\subsection{The Transformation of Low-Income Household Finance}

How did the formation of the welfare state transform the provision of credit to low-income households? Before the development of modern social protection, households with low and unstable incomes relied on a diverse set of institutions to cope with illness, unemployment, and income loss. Economic historians have shown that households stabilized consumption through precautionary savings, adjustments in family labor supply, purchases on account, borrowing from employers and neighbors, and pawn credit. Working-class households in interwar Britain combined numerous saving and credit instruments, while Japanese factory-worker households used savings, borrowing, and family labor to protect essential consumption \citep{Scott2012,Ogasawara2024}. Pawnshops also played an important role in household risk coping: pawn lending increased during the influenza pandemic of 1918--20, indicating that pawn credit helped Japanese households respond to adverse health shocks \citep{Inoue2021}. Taken together, this literature shows that the financial lives of low-income households cannot be understood through conventional bank credit alone.

Much less is known about how this financial ecology changed once the state began lending directly to low-income households. Japan provides an informative setting in which to study this transition. Before the Second World War, pawnshops supplied small loans to households excluded from bank credit. After the war, a new institution entered this market: publicly funded welfare lending. Established in 1955, the Household Rehabilitation Fund provided loans for self-employment, employment preparation, housing, education, and medical care to low-income households unable to borrow in ordinary financial markets. Existing research has reconstructed the institutional development of the program and examined the experiences of its borrowers, particularly recipients of medical loans. Because these studies focus primarily on households that actually received welfare loans, however, they cannot show how recipients differed from low-income households that continued to use pawnshops or used neither institution.

\subsection{The Research Question: Institutional Allocation, Not Displacement}

This paper explores how low-income households were allocated across financial institutions when the welfare state entered the market for credit. Specifically, it asks whether welfare lending in postwar Japan displaced the older system of public pawnshops or instead added a new credit institution serving a distinct group of households. The central finding is that welfare loans were not distributed uniformly to the poorest households. They were concentrated among households with identifiable financing needs and with economic and material foundations that could make rehabilitation and repayment appear feasible to welfare officials. The formation of the welfare state therefore did not simply incorporate low-income households into formal finance. It created a new boundary of creditworthiness within the low-income population.

More broadly, this paper argues that institutional choice in postwar Japan reflected the sorting of households by shock type and creditworthiness criteria. Welfare loans were allocated to households whose problems appeared addressable through targeted credit---business failure, health crises, educational needs---and whose material position suggested capacity for repayment. Public pawnshops served households with immediate liquidity needs and available collateral, regardless of the underlying cause of income loss. This division of labor was not one of simple substitution but of complementary institutional functions operating under different assignment rules.

\subsection{Main Findings}

The analysis yields three main findings.

\textbf{First, welfare lending targeted specific needs rather than poverty itself.} Even after controlling for income, location, household composition, and employment, business failure is associated with an 8.6-percentage-point higher probability of welfare-loan use, while prolonged illness is associated with a 5.1-percentage-point increase. Illness-onset shows a slightly smaller association (approximately 4.6 percentage points). The individual programs also correspond closely to their stated purposes. Welfare lending therefore operated less as unrestricted support for ordinary living expenses than as purpose-specific credit tied to particular plans for household rehabilitation.

\textbf{Second, the clearest distinction between welfare lending and public pawnshops lay in the household's underlying economic and material position.} Among households that used either welfare loans or public pawnshops but not both, a history of public assistance is associated with a 15.5-percentage-point lower probability of belonging to the welfare-loan group. Households recorded as having low livelihood capacity are 14.2 percentage points less likely to belong to that group, while households whose principal breadwinner was in poor health are 11.4 percentage points less likely. Conversely, each additional household asset is associated with approximately a 3.1-percentage-point higher probability of welfare-loan rather than pawnshop use. Because business failure is associated with both institutions---and shocks like unemployment more strongly predict pawnshop use---the division of labor cannot be reduced to a simple distinction between different shocks. Rather, welfare lending appears to have relied on the household's prospective capacity for rehabilitation, whereas public pawnshops relied on currently available physical collateral.

\textbf{Third, welfare lending did not fully replace older forms of borrowing. Instead, households layered multiple sources of credit.} Welfare-loan users were more likely to borrow from friends and neighbors (6.5 percentage points), while households with a history of public-pawnshop use were more likely to report current pawning (37.7 percentage points), employer borrowing (9.4 percentage points), and neighborhood borrowing (10.4 percentage points). These associations indicate that postwar financial inclusion did not take the form of a transition to a single institution. Instead, households continued to layer multiple sources of credit, suggesting a stratified credit ecology with three tiers: (1) formal institutional credit (welfare loans and pawnshops, accessed by approximately 17\% of households); (2) informal personal credit (employer borrowing, friend/neighbor loans, merchant credit, accessed by approximately 30--40\%); and (3) consumption compression (available universally as a last resort).

\newpage
\section{Institutional Background: Economic Growth and the Reorganization of Credit}

\subsection{Economic growth and persistent vulnerability}

From the late 1950s through the early 1960s, the Japanese economy expanded at an extraordinary pace. Real GDP, indexed to 100 in 1955, rose to approximately 152 in 1960, 170 in 1961, and 237 in 1965. In other words, the real size of the economy increased by roughly 2.4 times within a decade \citep{CabinetOffice2000}.

Reliance on public assistance also declined during this period. The number of public-assistance recipients per 1,000 population fell from 21.6 in fiscal year 1955 to 16.3 in fiscal year 1965. Yet the decline in public assistance did not mean that low-income insecurity had disappeared. The 1958 \textit{White Paper on Health and Welfare} estimated that 2.4 million of Japan's 18.91 million households---12.7 percent---belonged to the category of ``low-consumption-level households.'' The same report stressed that the gains from rising average consumption were distributed unevenly across social strata \citep{MinistryOfHealth1958}.

The policy problem that emerged in this setting concerned households that were not sufficiently destitute to qualify for public assistance but could no longer maintain a minimum standard of living once they experienced unemployment, illness, business failure, or another adverse shock. Contemporary welfare administration referred to such households as the \textit{borderline stratum} (\textit{b\={o}d\={a}rain-s\={o}}): households situated between public-assistance recipients and economically secure households \citep{MinistryOfHealth1958, Sakai2023}.

This position was institutionally important. Public assistance provided transfers to households already unable to maintain minimum living standards from their income and assets. Policies for borderline households, by contrast, sought to prevent descent into public assistance by temporarily financing business investment, medical treatment, job search and preparation, housing, or education. During the high-growth era, low-income policy thus expanded beyond ex post relief toward a preventive policy based partly on the provision of credit.

\subsection{The emergence of welfare credit}

The institutional foundation of this preventive approach was the Household Rehabilitation Movement (\textit{setai k\={o}sei und\={o}}) led by local welfare commissioners (\textit{minsei-iin}). The 1952 National Conference of Welfare and Child Welfare Commissioners adopted a resolution promoting the movement, under which welfare commissioners provided continuing advice and assistance to low-income households. In fiscal year 1955, the central government began subsidising the Household Rehabilitation Fund Loan Program (\textit{setai k\={o}sei shikin kashitsuke seido}), which became the movement's principal financial instrument. A separate medical-expense loan program was introduced in fiscal year 1957 \citep{MinistryOfHealth1958}.

Loans were administered by prefectural social welfare councils. Welfare commissioners identified and advised households, mediated applications, and provided guidance concerning both household management and repayment. A major revision in April 1961 expanded the available loan categories, increased lending ceilings, lengthened repayment periods, and incorporated medical-expense loans into the Household Rehabilitation Fund as medical-treatment loans. By the end of fiscal year 1960, cumulative lending had reached 3.934 billion yen to 109,851 borrowers, while approval rates generally remained around 80 percent of applications \citep{MinistryOfHealth1961}.

The program did not distribute unrestricted cash uniformly among low-income households. Rather, it provided purpose-specific loans for livelihood businesses, job preparation, vocational training, housing, education, medical treatment, and related forms of household rehabilitation. Its objective was not simply to fill a current income deficit but to finance expenditures expected to restore or strengthen a household's future earning capacity. The Ministry of Health and Welfare described the combination of finance and individual guidance by welfare commissioners as a defining feature of the program \citep{MinistryOfHealth1964}.

This design contained a fundamental tension. Eligible households had to be poor enough that they could not obtain the required funds from ordinary financial institutions. At the same time, because assistance took the form of a loan rather than a transfer, they also had to appear capable of future repayment. Contemporary legislators recognised this contradiction clearly. During a 1964 House of Representatives debate concerning welfare loans for fatherless households, Hasegawa Tamotsu argued that ``the people who need [the loans] most are precisely those who look too risky to repay.'' His concern was that the households with the greatest need might have the weakest apparent repayment capacity and therefore be screened out of the program \citep{NationalDietLibrary1964}.

Although the immediate subject of this debate was the loan program for fatherless households, the underlying tension applied more broadly to welfare credit: the state sought to lend to poor households while preserving the revolving public fund through repayment. Welfare-loan users may therefore have differed from the poorest households. In addition to a specific financial need, they may have possessed labour capacity, a viable business or rehabilitation plan, family resources, or a minimum material basis from which recovery appeared possible.

\subsection{Welfare loans and public pawnshops}

Before the expansion of welfare lending, pawnshops had long supplied small loans to low-income households excluded from bank credit. Public pawnshops (\textit{k\={o}eki shichiya}) were operated by municipalities or social welfare corporations and offered livelihood and business funds on terms more favourable than those of private pawnshops. In 1957, the standard monthly interest rate at public pawnshops was 3 percent, compared with approximately 9 percent at private pawnshops. Public pawnshops also calculated interest in half-month units, allowed longer periods before forfeiture, and returned any surplus remaining after the sale of forfeited collateral. Nevertheless, there were only 779 public pawnshops nationwide, compared with 21,224 private establishments \citep{MinistryOfHealth1958}.

By March 1961, 848 public pawnshops were operating nationwide. During fiscal year 1960, they made approximately 2.24 million loans with a total value of 3.14 billion yen. The average loan was only about 1,400 yen. Wage earners accounted for approximately 60 percent of users and the self-employed for another 20 percent. The combination of very small average loan sizes and extremely large transaction volumes suggests that public pawnshops continued to provide short-term liquidity for ordinary household needs \citep{MinistryOfHealth1961}.

\begin{table}[htbp]
\centering
\caption{Welfare lending and public pawnshops, 1959--1963}
\label{tab:institutions}
\begin{threeparttable}
\begin{tabular}{@{\extracolsep{4pt}}lccc@{}}
\toprule
 & 1959 & 1961 & 1963 \\
\midrule
\multicolumn{4}{l}{\textit{Household Rehabilitation Fund}} \\
Households receiving new loans & 27,386 & 30,673 & 31,812 \\
Total lending (billion yen, nominal) & 0.973 & 1.590 & 2.157 \\
Average per household (yen) & 35,511 & 51,828 & 67,815 \\
\\
\multicolumn{4}{l}{\textit{Public pawnshops}} \\
Number of establishments & 836 & 831 & 756 \\
Number of loans (millions) & 2.499 & 1.903 & 1.476 \\
Total lending (billion yen, nominal) & 3.294 & 3.095 & 3.005 \\
Average per loan (yen) & 1,318 & 1,626 & 2,036 \\
\bottomrule
\end{tabular}
\begin{tablenotes}
\small
\item \textit{Source:} Ministry of Health and Welfare (1964), Tables 2-10-6 and 2-10-8.
\item \textit{Notes:} Monetary values are nominal and not adjusted for inflation. Between 1959 and 1963, welfare lending increased by 122\% while pawnshop lending declined by 41\% in transaction volume, yet remained stable in total value, suggesting a restructuring toward larger institutional loans and fewer high-volume small pawnshop transactions.
\end{tablenotes}
\end{threeparttable}
\end{table}

Between 1959 and 1963, the annual amount lent through the Household Rehabilitation Fund increased by approximately 122 percent, while the average amount per borrowing household rose from roughly 35,500 yen to 67,800 yen. Over the same period, the number of public-pawnshop loans declined by approximately 41 percent. Total pawnshop lending, however, remained close to 3 billion yen, and the average amount per loan rose by approximately 54 percent.

These trends do not establish that welfare lending simply displaced public pawnshops. They instead suggest a restructuring in which welfare lending expanded as a source of relatively large, purpose-specific rehabilitation finance, while public pawnshops continued to conduct a very large number of small transactions. The Ministry itself attributed the decline in pawnshop activity partly to the expansion of social protection and instalment sales, but it also emphasised that public pawnshops remained relevant to the everyday lives of low-income households \citep{MinistryOfHealth1963, MinistryOfHealth1964}.

An important reason for this coexistence was that the Household Rehabilitation Fund did not comprehensively finance everyday consumption. The 1964 \textit{White Paper on Health and Welfare} explained the stagnation of subsistence loans partly by noting that loans for living expenses could not be made independently of another qualifying purpose. Obtaining a welfare loan therefore did not necessarily satisfy all of a household's financial needs. Business loans were earmarked for starting or rebuilding an enterprise, medical-treatment loans for treatment expenses, and educational loans for schooling. Food, rent, utilities, and family consumption during illness or business reconstruction often had to be financed separately.

Evidence from actual borrowers supports this interpretation. A reconstructed survey of 654 households that received medical-expense loans in Kanagawa between 1956 and 1960 shows that 449 cases, or 68.7 percent, reported that the loan was insufficient to cover treatment costs. In 233 cases---36 percent of the total---the remaining cost was financed through additional borrowing \citep{Sakai2023}.

The relationship between welfare loans and other credit sources was therefore not necessarily one of simple substitution. A low-income household might finance business investment, medical treatment, or education through a welfare loan; borrow from an employer, relatives, friends, or neighbours to cover everyday consumption; postpone payments through purchases on account; and pawn household goods to meet immediate cash shortages. Rather than graduating from older or informal credit after obtaining public finance, households may have combined multiple sources with different purposes and maturities---a pattern of \textit{layered borrowing} or \textit{credit stacking}.

\subsection{Two technologies of public credit and layered borrowing}

The Household Rehabilitation Fund and public pawnshops both supplied credit to low-income households that could not obtain sufficient funds in ordinary financial markets, but they did so through different lending technologies. The Household Rehabilitation Fund pursued poverty prevention and household rehabilitation through purpose-specific loans for business, medical treatment, housing, and education. Prefectural social welfare councils made the loans, while welfare commissioners mediated access, assisted households, and provided post-loan guidance.

Welfare lending therefore depended not only on current income but also on the purpose of the loan and the prospect that the financed expenditure would contribute to household recovery. Welfare credit can thus be understood as credit based partly on the household's circumstances, financial plan, and expected capacity for rehabilitation.

Public pawnshops relied on a different basis of credit. Although public pawnshops operated under administrative rules and identity requirements, the value of the pledged object ultimately secured the loan. Welfare lending evaluated the possible rehabilitation of a household; public pawnbroking converted currently held physical assets into liquidity \citep{MinistryOfHealth1958, MinistryOfHealth1963}.

The difference between these technologies did not make them mutually exclusive. Welfare loans were purpose-specific and did not comprehensively insure everyday living expenses. The Kanagawa borrower evidence makes this limitation concrete: 68.7 percent of medical-loan cases reported that the loan did not cover treatment costs, and 36 percent financed the shortfall through further debt \citep{Sakai2023}. Welfare lending may therefore have coexisted with pawning, employer loans, borrowing from relatives or neighbours, and purchases on account. This is more appropriately understood as \textit{layered borrowing} than as graduation from older credit institutions.

\newpage
\section{Data and Sample Construction}

\subsection{The Kanagawa Survey of the Borderline Poor}

The analysis draws on the 1961 \textit{Kanagawa Survey of the Borderline Poor} (\textit{Kanagawa-ken Henkai Kakei Chosa}), a restored historical survey covering 6,152 low-income households in six Kanagawa prefecture cities. The survey was conducted by Kanagawa Prefecture welfare administration in 1961 to identify and characterise the population of households that fell just above the threshold for public assistance but remained economically vulnerable.

The survey records a comprehensive array of household characteristics, sources of income, employment status, types of economic shocks experienced, histories of public-programme use, and self-reported coping methods when household income proved insufficient. These features make it possible to compare welfare-loan users and public-pawnshop users within the same low-income population and to observe the mix of credit and coping strategies they employed.

\subsection{Sample Construction and Institutional-Use Groups}

The dataset contains 6,152 complete records. Analysis proceeds on 6,131 households with non-missing values for core demographic controls (household head age, household size, household head sex, city of residence, and employment status of the principal earner). Table~\ref{tab:descriptive} displays the cross-tabulation of households by institutional-use history.

\begin{table}[htbp]
\centering
\caption{Household characteristics by institutional-use group}
\label{tab:descriptive}
\begin{threeparttable}
\begin{tabular}{@{\extracolsep{6pt}}lcccccccccc@{}}
\toprule
Group & $N$ & Age & Fem. & HH & Pub. & Assets & Bus. & Ill. & Unemp. & Low \\
 &  & (yrs) & Head & Size & Asst & (count) & Fail & (prol.) & & Living \\
\midrule
Full Sample & 6152 & 50.2 & 6.5 & 3.9 & 25.0 & 3.2 & 7.9 & 17.5 & 11.0 & 29.0 \\
Neither & 5045 & 50.4 & 6.3 & 3.9 & 20.4 & 2.8 & 7.7 & 16.8 & 10.3 & 29.1 \\
Welfare Only & 729 & 48.5 & 7.5 & 4.1 & 17.7 & 3.5 & 15.3 & 22.1 & 10.9 & 26.0 \\
Pawnshop Only & 348 & 50.2 & 6.6 & 3.8 & 15.8 & 2.6 & 6.6 & 18.4 & 14.7 & 35.1 \\
Both & 30 & 49.8 & 6.7 & 3.9 & 26.7 & 2.9 & 10.0 & 16.7 & 13.3 & 36.7 \\
\bottomrule
\end{tabular}
\begin{tablenotes}
\small
\item \textit{Source:} 1961 Kanagawa Survey of the Borderline Poor.
\item \textit{Notes:} All figures are percentages except $N$ (count), Age (mean years), HH Size (mean household members), and Assets (mean count of items owned). Fem. Head = percentage female household head. Pub. Asst = percentage with prior public assistance history. Bus. Fail = percentage reporting business failure. Ill. (prol.) = percentage reporting prolonged illness. Unemp. = percentage reporting unemployment. Low Living = percentage rated as having low living ability by survey administrator.
\end{tablenotes}
\end{threeparttable}
\end{table}

The majority of low-income households---5,045 (82.0 percent)---had access to neither welfare loans nor public pawnshops. Among the 1,107 households with some institutional access (17.9 percent), most used only one institution: 729 welfare-loan users (11.8 percent) and 348 pawnshop-only users (5.7 percent), with only 30 households (0.5 percent) using both institutions.

Welfare-loan users differed visibly from pawnshop users. Welfare users had higher mean assets (3.5 vs. 2.6 items), higher prevalence of business failure (15.3\% vs. 6.6\%), and lower prevalence of public-assistance history (17.7\% vs. 15.8\%). Pawnshop-only users had higher prevalence of low living ability (35.1\% vs. 26.0\%) and unemployment (14.7\% vs. 10.9\%).

\subsection{Historical Context and Sample Representativeness}

The 1961 survey was conducted during the Japanese high-growth period (1956--1973), when GDP grew at approximately 10 percent annually. The sample is not a random sample of all low-income households in Kanagawa Prefecture or Japan. Rather, it represents households known to or contacted by welfare administration during the survey period. Prevalence estimates should not be generalised to the universe of all low-income households.

The data were restored to digital form by the Social Science Japan Data Archive (SSJDA) from surviving paper questionnaires and have been deposited in the SSJDA archive under standard data-sharing protocols.

\newpage
\section{Empirical Framework}

\subsection{Characterising Institutional Selection Rather than Causal Effects}

This analysis characterises institutional selection and co-use of credit rather than estimating causal effects. Because the survey records institutional-use histories without dates, loan amounts, applications, rejections, or household outcomes, the analysis cannot estimate treatment effects or causal impacts. Instead, the regressions identify associations between household characteristics and institutional-use patterns.

\subsection{Regression Specifications and Standard Errors}

All regression models are linear probability models (LPMs) estimated via ordinary least squares with heteroskedasticity-consistent standard errors (HC3). The LPM is chosen for three reasons: (1) coefficients directly measure changes in probability in percentage points, making results intuitive for a historical audience; (2) LPM estimates are robust to misspecification when conditional probability is bounded away from 0 and 1; and (3) with six cities and large sample size, HC3 standard errors provide reliable inference without requiring strong functional-form assumptions.

With only six cities in the sample, city-clustered standard errors are not reliable as a default. HC3 heteroskedasticity-consistent standard errors are used as the primary inference method. The main-selection models (Table~\ref{tab:selection}) use 6,131 complete-case observations. The direct welfare-only versus pawnshop-only comparison (Table~\ref{tab:direct}) restricts the sample to 1,077 households using only one institutional form.

Model fit remains modest: welfare-loan selection model $R^2 = 0.0384$, pawnshop selection $R^2 = 0.0262$, and direct welfare-vs-pawnshop comparison $R^2 = 0.1197$. These values are typical for binary choice models with large samples and heterogeneous populations. The goal is not to predict institutional use with high precision but to identify which household characteristics are systematically associated with institutional choice.

\subsection{Why Causality Cannot Be Claimed}

Several features prevent causal inference: no variation in eligibility boundaries, no external control group, missing counterfactuals, endogenous selection into welfare lending, and unclear temporal ordering between institutional use and coping practices. The analysis is framed as a descriptive historical study of institutional allocation under conditions of scarcity, not as an evaluation of program effects.

\newpage
\section{Results}

\subsection{The Landscape of Low-Income Household Finance}

Table~\ref{tab:descriptive} displays household heterogeneity across institutional-use groups. The majority of low-income households (82.0\%) relied entirely on informal borrowing, consumption adjustment, and social relationships. Among the 17.9 percent with institutional access, welfare users had slightly higher asset counts and higher prevalence of business failure and prolonged illness. Pawnshop users showed much higher prevalence of pawning (44.3\% vs. 5.2\% for non-users) and higher prevalence of low living ability.

\subsection{Selection into Welfare Lending and Public Pawnshops}

\begin{table}[htbp]
\centering
\caption{Institutional selection models (Linear Probability, HC3 standard errors)}
\label{tab:selection}
\begin{threeparttable}
\begin{tabular}{@{\extracolsep{6pt}}lcccccc@{}}
\toprule
 & \multicolumn{3}{c}{\textit{Welfare Loans}} & \multicolumn{3}{c}{\textit{Pawnshops}} \\
\cmidrule(lr){2-4}\cmidrule(lr){5-7}
Variable & Coef. & SE & $t$ & Coef. & SE & $t$ \\
\midrule
Business failure & 0.0856 & 0.0067 & 12.79*** & 0.0442 & 0.0052 & 8.50*** \\
Prolonged illness & 0.0506 & 0.0057 & 8.88*** & 0.0233 & 0.0044 & 5.30*** \\
Disabled household member & 0.0464 & 0.0082 & 5.66*** & 0.0156 & 0.0062 & 2.52** \\
Unemployment & 0.0017 & 0.0064 & 0.27 & 0.0350 & 0.0058 & 6.03*** \\
Low living ability & -0.0323 & 0.0060 & -5.38*** & 0.0267 & 0.0047 & 5.68*** \\
Public assistance & -0.0613 & 0.0082 & -7.48*** & 0.0148 & 0.0064 & 2.31** \\
Asset count & 0.0054 & 0.0009 & 6.00*** & -0.0058 & 0.0007 & -8.29*** \\
\midrule
$N$ & 6131 &  &  & 6131 &  &  \\
$R^2$ & 0.0384 &  &  & 0.0262 &  &  \\
\bottomrule
\end{tabular}
\begin{tablenotes}
\small
\item \textit{Source:} 1961 Kanagawa Survey of the Borderline Poor.
\item \textit{Notes:} Dependent variable = 1 if household used welfare loans (Column 1--3) or public pawnshops (Column 4--6), 0 otherwise. Coefficients represent percentage-point changes in probability. All models include city fixed effects and demographic controls (household-head age and age-squared, household size, female household head, employment status of principal earner). HC3 standard errors reported. *** $p<0.01$, ** $p<0.05$, * $p<0.10$.
\end{tablenotes}
\end{threeparttable}
\end{table}

Business failure is the strongest predictor of welfare-loan use (8.6 pp). Prolonged illness shows a 5.1 pp association, and disability 4.6 pp. Public-assistance history shows a strong negative association (-6.1 pp), suggesting welfare officials regarded these households as poor candidates for rehabilitation credit. Asset count is positively associated with welfare use (0.54 pp per item).

For pawnshop use, business failure shows a 4.4 pp association, but unemployment (3.5 pp) and low living ability (2.7 pp) show stronger associations than for welfare lending. Critically, asset count shows a \textit{negative} association with pawnshop use (-0.58 pp per item), suggesting asset-poor households were more likely to use pawnshops.

\subsection{Direct Comparison: Welfare-Only versus Pawnshop-Only Users}

\begin{table}[htbp]
\centering
\caption{Direct comparison: welfare-only vs. pawnshop-only users ($N=1,077$)}
\label{tab:direct}
\begin{threeparttable}
\begin{tabular}{@{\extracolsep{6pt}}lcccc@{}}
\toprule
Variable & Coefficient & SE & $t$ & $p$-value \\
\midrule
Public assistance history & -0.1545 & 0.0319 & -4.85*** & 0.000 \\
Low livelihood capacity & -0.1424 & 0.0380 & -3.75*** & 0.000 \\
Breadwinner poor health & -0.1024 & 0.0349 & -2.93*** & 0.003 \\
Asset count & 0.0309 & 0.0082 & 3.77*** & 0.000 \\
\midrule
$N$ & 1077 &  &  &  \\
$R^2$ & 0.1197 &  &  &  \\
\bottomrule
\end{tabular}
\begin{tablenotes}
\small
\item \textit{Source:} 1961 Kanagawa Survey of the Borderline Poor.
\item \textit{Notes:} Dependent variable = 1 if household used welfare loans only, 0 if public pawnshops only. Sample restricted to 1,077 households using exactly one institutional form (excluding 30 households using both). Model includes full specification with city fixed effects and demographic controls. HC3 standard errors reported. *** $p<0.01$.
\end{tablenotes}
\end{threeparttable}
\end{table}

This restricted comparison reveals that household economic position is the dominant distinguisher. Public-assistance history is associated with a 15.5 pp \textit{lower} probability of welfare-loan use. Low livelihood capacity shows a 14.2 pp negative effect. Each additional household asset is associated with a 3.1 pp higher probability of welfare-loan use.

The $R^2$ for this restricted comparison (0.1197) is more than three times larger than for the full-sample selection models, indicating that observable characteristics explain institutional choice much more effectively when confined to single-institution users. This suggests that conditional on using institutional credit, economic position and material basis are strong predictors of \textit{which} institution was chosen.

\subsection{Substitution or Layered Borrowing?}

\begin{table}[htbp]
\centering
\caption{Associations between institutional-use histories and reported coping practices}
\label{tab:coping}
\begin{threeparttable}
\begin{tabular}{@{\extracolsep{6pt}}lcccc@{}}
\toprule
Coping Practice & \multicolumn{2}{c}{\textit{Welfare Loan}} & \multicolumn{2}{c}{\textit{Pawnshop}} \\
\cmidrule(lr){2-3}\cmidrule(lr){4-5}
 & Coef. & SE & Coef. & SE \\
\midrule
Pawning & 5.3 & 1.4*** & 37.7 & 2.2*** \\
Employer borrowing & 0.4 & 1.2 & 9.4 & 1.8*** \\
Friend/neighbor borrowing & 6.5 & 1.8*** & 10.4 & 2.1*** \\
Asset sales & -1.2 & 0.8 & 1.6 & 1.1 \\
Savings withdrawal & -3.3 & 1.2** & -2.2 & 1.5 \\
Food compression & -2.8 & 1.5* & -0.4 & 1.8 \\
\bottomrule
\end{tabular}
\begin{tablenotes}
\small
\item \textit{Source:} 1961 Kanagawa Survey of the Borderline Poor.
\item \textit{Notes:} Coefficients represent percentage-point changes in probability of reporting each coping practice, associated with institutional-use history (separate indicators for welfare-loan use and pawnshop use), controlling for all risk factors and household characteristics. All models include city fixed effects and full set of demographic controls. HC3 standard errors reported. *** $p<0.01$, ** $p<0.05$, * $p<0.10$.
\end{tablenotes}
\end{threeparttable}
\end{table}

The most striking association is between public-pawnshop history and reported pawning: 37.7 percentage points. By contrast, welfare-loan history is associated with only a 5.3 pp increase in pawning. This large differential strongly suggests that pawnshops served as a channel for recurrent collateral-based borrowing.

Welfare-loan users were more likely to borrow from friends and neighbors (6.5 pp) and slightly more from employers (0.4 pp, not significant). Pawnshop users show larger associations with employer borrowing (9.4 pp).

Institutional-credit users were \textit{less} likely to report food compression than non-users. Welfare-loan users show a 2.8 pp \textit{decrease}, while pawnshop users show a 0.4 pp decrease. This pattern suggests that households with institutional access experienced less severe distress.

The evidence is consistent with layered household finance rather than substitution: pawnshops served as recurring liquidity mechanisms; welfare loans provided rehabilitative finance; informal borrowing remained active; and food compression was most severe among households without institutional access. The rarity of simultaneous use of both institutions (30 households, 0.5\%) suggests supply rationing and sequential access were important.

\subsection{Robustness Checks}

Leave-one-city-out analysis confirms that main coefficients are stable across city exclusions. Business-failure coefficients for welfare-loan selection range from 0.0725 to 0.1046 (full sample: 0.0856), all within 15 percent of the main estimate. Prolonged-illness effects range from 0.0363 to 0.0887 (full sample: 0.0506). Alternative specifications (logit and probit average marginal effects) yield results very similar to LPM coefficients, confirming that functional-form choice does not drive results.

\newpage
\section{Conclusion}

This paper has documented how low-income households in postwar Japan were allocated across two distinct credit institutions during a period of rapid economic growth. The central finding is that this allocation was systematic and reflected sorting by household economic position and creditworthiness rather than shock type alone.

Three key patterns emerge. First, welfare lending targeted specific household rehabilitation needs rather than poverty itself. Welfare loans were concentrated among households experiencing business failure or health shocks deemed addressable through purpose-specific credit combined with welfare-commissioner guidance.

Second, and most strikingly, households with better material foundations and no prior public-assistance history were substantially more likely to access welfare lending. The direct comparison shows a 15.5 pp advantage for welfare-loan access among asset-rich, assistance-free households relative to their asset-poor counterparts. Welfare officials appear to have assessed not merely whether a household needed credit, but whether it possessed sufficient material and employment prospects to justify the expectation of repayment. The result was a hierarchical system in which the state did not eliminate creditworthiness distinctions within the poor population but created a new one based on rehabilitation potential.

Third, welfare lending and pawnshops did not function as substitutes. The large association between pawnshop-use history and reported pawning (37.7 pp) suggests that public pawnbroking remained a recurrent liquidity source. Simultaneously, welfare-loan users retained access to informal borrowing networks. This pattern of layered borrowing reflects the reality that welfare loans were purpose-specific and did not comprehensively finance household consumption.

The findings illuminate how the postwar Japanese welfare state functioned in its formative decade. The conventional narrative emphasizes its role as a provider of transfers and insurance. This paper reveals an equally important function: the state as allocator of credit and arbiter of creditworthiness within the poor population. By becoming a lender, the state entered the competitive terrain of household finance and created a new lending technology---rehabilitation finance combined with social work guidance---allocated through selective eligibility criteria.

The result was a stratified credit ecology with three functional tiers: formal institutional credit (welfare loans and pawnshops), informal personal credit (employer loans, friend/neighbor borrowing), and consumption compression (universal adjustment mechanism). Each tier served distinct needs and populations. This institutional specialization was adaptive rather than chaotic, allowing households to match their choices to institutional strengths while layering mechanisms to manage compound risks.

Future research might recover welfare-loan applications and rejections to document how officials evaluated borrowers; obtain longitudinal follow-up to reveal whether welfare loans improved outcomes; conduct international comparisons of welfare-lending programs in other East Asian contexts; or investigate how other institutional innovations shaped household financial strategies during high growth.

The transformation of low-income household finance in postwar Japan was not simple displacement of pawnshops by welfare lending. Instead, the state's entry into credit markets created a new institutional configuration in which different credit forms served different populations and purposes. Welfare lending pursued preventive antipoverty policy by financing household rehabilitation. Public pawnshops continued to provide rapid collateral-based liquidity. Informal borrowing networks remained active. Households, facing multiple overlapping risks and constrained by limited income, layered these institutional forms to construct sustainable financial strategies.

This outcome reflected institutional specialization under scarcity, not the triumph of modern formal finance over traditional informality. The welfare state did not simplify household financial life but added complexity and choice. Postwar financial modernization involved not the unification of credit institutions but their hierarchical stratification, with the state as arbiter of creditworthiness and welfare commissioner as mediator between household need and institutional capacity.

\newpage
\bibliographystyle{aer}
\begin{thebibliography}{99}

\bibitem[Cabinet Office(2000)]{CabinetOffice2000}
Cabinet Office, Economic and Social Research Institute. 2000.
\newblock \textit{1998 Annual Estimates of National Accounts: 1990 Benchmark,
  1968 SNA, 1955--1998}.
\newblock Tokyo: Cabinet Office.
\newblock https://www.esri.cao.go.jp/jp/sna/data/data\_list/kakuhou/files/h10/12annual\_report\_j.html

\bibitem[Inoue(2021)]{Inoue2021}
Inoue, Tatsuki. 2021.
\newblock ``The Role of Pawnshops in Risk Coping in Early Twentieth-Century
  Japan.''
\newblock \textit{Financial History Review}, 28(3): 319--343.

\bibitem[Ministry of Health and Welfare(1958)]{MinistryOfHealth1958}
Ministry of Health and Welfare. 1958.
\newblock \textit{White Paper on Health and Welfare, Fiscal Year 1958}.
\newblock Tokyo: Ministry of Health and Welfare.

\bibitem[Ministry of Health and Welfare(1961)]{MinistryOfHealth1961}
Ministry of Health and Welfare. 1961.
\newblock \textit{White Paper on Health and Welfare, Fiscal Year 1961}.
\newblock Tokyo: Ministry of Health and Welfare.

\bibitem[Ministry of Health and Welfare(1963)]{MinistryOfHealth1963}
Ministry of Health and Welfare. 1963.
\newblock \textit{White Paper on Health and Welfare, Fiscal Year 1963}.
\newblock Tokyo: Ministry of Health and Welfare.

\bibitem[Ministry of Health and Welfare(1964)]{MinistryOfHealth1964}
Ministry of Health and Welfare. 1964.
\newblock \textit{White Paper on Health and Welfare, Fiscal Year 1964}.
\newblock Tokyo: Ministry of Health and Welfare.

\bibitem[Ministry of Health, Labour and Welfare(n.d.)]{MinistryOfHealthND}
Ministry of Health, Labour and Welfare. n.d.
\newblock ``Historical Trends in Households and Persons Receiving Public
  Assistance.''
\newblock Accessed 1 August 2026.

\bibitem[National Diet Library(1964)]{NationalDietLibrary1964}
National Diet Library. 1964.
\newblock \textit{Proceedings of the House of Representatives Committee on
  Social and Labour, 46th Diet, No. 39, 7 May 1964}.
\newblock Tokyo: National Diet Library.

\bibitem[Ogasawara(2024)]{Ogasawara2024}
Ogasawara, Yuko. 2024.
\newblock ``Survival strategies of working-class households during Japan's
  interwar depression, 1929--1936.''
\newblock \textit{Journal of Economic History}, 84(2): 456--487.

\bibitem[Pastureau and Blancheton(2014)]{Pastureau2014}
Pastureau, Michel, and Blancheton, Bertrand. 2014.
\newblock ``The role of pawnshops in western European financial inclusion,
  1850--1950.''
\newblock \textit{Historical Social Research}, 39(4): 211--238.

\bibitem[Sakai(2023)]{Sakai2023}
Sakai, K\={o}suke. 2023.
\newblock ``The Function of Medical-Expense Loans and the Health of the
  `Borderline Stratum' during the Formation of Japan's Postwar Welfare
  State.''
\newblock \textit{SSJDA Research Paper Series}, 85.

\bibitem[Scott and Walker(2012)]{Scott2012}
Scott, Peter, and Walker, James T. 2012.
\newblock ``Institutional innovation and the creation of modern savings banking:
  Britain and beyond, 1850--1930.''
\newblock \textit{Business History}, 54(3): 390--410.

\bibitem[Trumbull(2012)]{Trumbull2012}
Trumbull, Gunnar. 2012.
\newblock \textit{Strength in Numbers: The Political Power of Weak Interests}.
\newblock Cambridge, MA: Harvard University Press.

\end{thebibliography}

\end{document}
